\section{Estrategia Empírica}

\subsection{Datos}

\paragraph{Fuentes}

La investigación combina dos fuentes, todas de acceso público. La primera es el registro histórico de precios de la Comisión Nacional de Energía (CNE). Cada observación corresponde a una actualización de precio informada por una EDS para un combustible determinado, con su fecha, y viene acompañada de la identificación de la estación, su distribuidor, sus coordenadas geográficas y su comuna y región. Como las EDS están obligadas por norma a informar cada variación de precio, el registro constituye un censo de la evolución de precios del mercado y no una muestra. La serie utilizada cubre desde enero de 2012 hasta julio de 2026.\footnote{Adicionalmente la CNE dispone información de servicios complementarios que ofrece cada EDS: tienda de conveniencia, lavado, baño público, farmacia, surtidor de camiones, entre otros. Estos atributos se emplean únicamente con fines descriptivos y de balance, ya que solo se dispone de la última actualización de cada estación y no de su evolución en el tiempo.}

La segunda es la serie semanal del precio mayorista con MEPCO, publicada por el Ministerio de Hacienda. Cubre desde el 7 de agosto de 2014 hasta el 2 de julio de 2026 y consta de 622 semanas consecutivas sin interrupciones, todas con vigencia a partir del día jueves, tal como establece la ley \parencite{LEY20765}. La serie está disponible para las gasolinas de 93 y 97 octanos y para el diésel, no existe una referencia MEPCO para la gasolina de 95 octanos, que por ello queda fuera de todo análisis de márgenes.

\paragraph{Construcción del panel}

El registro de la CNE genera una fila cada vez que una estación informa un precio, de modo que las observaciones son irregulares en el tiempo y no constituyen un panel. Su transformación en un panel balanceado descansa en tres decisiones, cuyo detalle e implementación se documentan en el Anexo~\ref{anexo:datos}. La unidad de análisis es la estación física y no el identificador administrativo, porque un mismo local puede aparecer en el registro bajo más de un código. Los identificadores que corresponden a un mismo local se fusionan en una sola unidad, de no hacerlo, un cambio de operador se contaría como una salida seguida de una entrada, y un registro duplicado inflaría el conteo de competidores cercanos.

El panel se construye a frecuencia semanal sobre la grilla de jueves del MEPCO, y de ahí se colapsa a frecuencia mensual. La semana no es una elección arbitraria, es la cadencia efectiva con que las estaciones reportan y es también la frecuencia nativa del precio mayorista. En esa grilla, el $81{,}1\%$ de las celdas corresponde a un precio efectivamente informado esa semana y el resto a un precio vigente arrastrado desde el último reporte, con un tope que impide arrastrarlo a través de períodos en que la estación no operaba. El margen bruto se define como el precio de venta menos el precio mayorista con MEPCO vigente esa semana, conforme a la \Cref{prop:margen}. Al ser el MEPCO una referencia nacional, el costo se adjudica por combustible y semana, no por estación.

El resultado es un panel semanal de 4.663.553 observaciones sobre 2.030 estaciones y 758 semanas, y su colapso mensual de 283.695 observaciones sobre 175 meses.

\paragraph{Definición de los eventos y de los grupos de comparación} Una entrada se define como el primer mes en que una estación física aparece en el registro, sujeto a tres restricciones que determinan qué eventos pueden usarse como tratamiento y que el Anexo~\ref{anexo:datos} detalla. Con estas restricciones, una incumbente se clasifica como tratada si experimenta una entrada dentro de un radio de 2 kilómetros mientras ya está operando. Las restantes conforman el grupo de comparación y se subdividen en un anillo, entre 2 y 5 kilómetros del evento más cercano, y un grupo lejano, a más de 5 kilómetros de toda entrada. La \Cref{tab:resumen} resume la composición resultante de la muestra.\footnote{Las salidas requieren una advertencia. El cambio de régimen de reporte de la CNE entre los archivos antiguos y los actuales deja una cicatriz en los datos que impide fechar con confianza el cierre de una estación en torno al quiebre, el Anexo~\ref{anexo:datos} documenta su origen y su magnitud. La limitación no afecta al diseño, que se construye exclusivamente sobre entradas.}

\begin{table}[H]
\centering
\caption{Composición de la muestra}
\label{tab:resumen}
\tabla{tab_desc_resumen}
\end{table}

La tabla y las figuras que siguen tienen por objeto mostrar que el panel reúne las condiciones que el diseño de identificación requiere, tales como número suficiente de eventos, dispersos en el tiempo, con grupos de comparación bien poblados y con variables de resultado que se comportan de manera estable.

La \Cref{fig:entradas} presenta las entradas por año y por marca de la estación entrante, desde 2014, que es el período en que rige el MEPCO y del que provienen las cohortes de tratamiento. Los eventos se distribuyen a lo largo de toda la serie. Destaca que la bandera blanca es el principal entrante del mercado, con 184 de las 482 entradas, por sobre Copec con 140.

\begin{figure}[H]
\centering
\includegraphics[width=\textwidth]{desc_entradas_por_anio.pdf}
\caption{Entradas por año y marca de la estación entrante}
\label{fig:entradas}
\end{figure}

La \Cref{fig:precio} muestra la evolución del precio promedio nacional por combustible. Las cuatro series se mueven juntas y están dominadas por el costo mayorista, que es común a todas las estaciones: es precisamente esa componente la que el margen descuenta. La \Cref{fig:margen} presenta el margen bruto promedio a partir de agosto de 2014. A diferencia del precio, el margen se mueve en un rango acotado, lo que confirma que descontar el MEPCO aísla una variable de nivel comparable a lo largo del tiempo y no una tendencia macroeconómica.

\begin{figure}[H]
\centering
\includegraphics[width=0.9\textwidth]{desc_precio_promedio.pdf}
\caption{Precio promedio nacional por combustible}
\label{fig:precio}
\end{figure}

\begin{figure}[H]
\centering
\includegraphics[width=0.9\textwidth]{desc_margen_promedio.pdf}
\caption{Margen bruto promedio sobre el precio mayorista con MEPCO}
\label{fig:margen}
\end{figure}

Por último, la \Cref{fig:mercado} describe el mercado local en el momento en que ocurre cada entrada. El $72{,}6\%$ de las entradas tiene al menos una estación vecina dentro de 2 kilómetros, con un promedio de 3,4 vecinos y una mediana de 2. Dos rasgos importan para la interpretación. Primero, los mercados afectados son efectivamente locales y poblados, de modo que la entrada constituye un cambio apreciable en la estructura competitiva y no una perturbación marginal. Segundo, solo el $19\%$ de las estaciones vecinas comparte marca con la entrante, lo que indica que la entrada típica es la llegada de un competidor y no la expansión de la red de una misma cadena.

\begin{figure}[H]
\centering
\includegraphics[width=\textwidth]{desc_mercado_local.pdf}
\caption{Composición de marca del mercado local al momento de la entrada}
\label{fig:mercado}
\end{figure}

\subsection{Identificación}

La estrategia sigue de cerca a \textcite{FISCHER2025}, cuyo objeto de estudio el efecto de la entrada de una estación sobre los precios de las incumbentes de su mercado local es el mismo de este trabajo. Se estima un estudio de eventos sobre el panel mensual de estaciones,
%
\begin{equation}
    y_{it} = \alpha_i + \lambda_t + \gamma_{r(i)a(t)} + \delta_{m(i)a(t)}
             + \sum_{\substack{\tau=-4 \\ \tau \neq -1}}^{4}
               \beta_\tau \, D^{\tau}_{it} + \varepsilon_{it},
    \label{eq:es}
\end{equation}
%
donde $y_{it}$ es el precio en logaritmo o el margen de la estación $i$ en el mes $t$, $\alpha_i$ es un efecto fijo de estación; $\lambda_t$ uno de mes, $\gamma_{r(i)a(t)}$ uno de región por año y $\delta_{m(i)a(t)}$ uno de marca por año. El indicador $D^{\tau}_{it}$ vale uno si la estación $i$ es tratada y el mes $t$ cae en el $\tau$-ésimo bin de seis meses respecto de la entrada que la trata. Los bins se agrupan en los extremos, de modo que $\tau \le -4$ y $\tau \ge 4$ acumulan todo lo que ocurre más allá de la ventana, y $\tau = -1$ es la categoría omitida. La ventana de cuatro bins previos y cinco posteriores, así como el agrupamiento de los extremos, replican el formato de reporte de \textcite{FISCHER2025}. Los errores estándar se agrupan a nivel de comuna.

El efecto promedio se obtiene de la versión estática de \eqref{eq:es},
%
\begin{equation}
    y_{it} = \alpha_i + \lambda_t + \gamma_{r(i)a(t)} + \delta_{m(i)a(t)}
             + \beta \, \mathrm{Post}_{it} + \varepsilon_{it},
    \label{eq:did}
\end{equation}
%
donde $\mathrm{Post}_{it}$ vale uno para una estación tratada desde el mes de la entrada en adelante.

Ambas especificaciones se estiman por efectos fijos bidireccionales sobre un diseño de adopción escalonada, situación en la que el estimador es un promedio ponderado de comparaciones entre cohortes que puede recibir ponderaciones negativas cuando el efecto del tratamiento varía entre ellas \parencite{CALLAWAYSANTANNA2021, SUNABRAHAM2021, DECHAISEMARTIN2020, BORUSYAK2024}. El Anexo~\ref{anexo:robustez} reestima el efecto promedio con el estimador de \textcite{SUNABRAHAM2021}, que nunca emplea como control a una incumbente ya tratada, y obtiene el mismo resultado; es la misma verificación que realizan \textcite{FISCHER2025} sobre su propio diseño.

Dos elementos de \eqref{eq:es} merecen justificación. El efecto fijo de estación absorbe todo atributo invariante del local tales como su ubicación, tamaño, y posición relativa dentro del mercado local, de modo que la identificación proviene únicament  de la variación dentro de cada estación. El efecto fijo de marca por año, que \textcite{FISCHER2025} no incorpora, responde a un rasgo propio del caso
chileno documentado en la sección anterior.\footnote{Bajo el modelo de consignación el precio de venta al público lo determina la compañía mayorista, y con tres cadenas que concentran cerca del $80\%$ de las EDS, un movimiento nacional de cualquiera de ellas es un confusor que ni el efecto fijo de estación ni el de región por año absorben.}

El supuesto de identificación es que el momento en que la entrada se materializa es exógeno desde la perspectiva de la incumbente, condicional en los efectos fijos de \eqref{eq:es}. Es el supuesto de \textcite{ARCIDIACONO2020}, adoptado también por \textcite{FISCHER2025} y por \textcite{KOH2022}. No exige que la ubicación de la entrada sea exógena (de hecho, no lo es) sino solo que, dado que una entrada va a ocurrir en un mercado, la fecha precisa  en que ocurre no responda a la evolución de los precios de las incumbentes. Este supuesto se defiende por las regulaciones que conlleva abrir una nueva estación y se justifica empíricamente más adelante.

\textcite{FISCHER2025} comparan las incumbentes tratadas contra el conjunto de estaciones no tratadas, sin construir un anillo de control. Este trabajo adopta el mismo criterio, y la \Cref{tab:control} justifica la elección estimando \eqref{eq:did} bajo las tres definiciones posibles del grupo de comparación: el anillo de estaciones entre 2 y 5 kilómetros del evento, las estaciones lejanas a más de 5 kilómetros de toda entrada, y la unión de ambas, que es el conjunto de las nunca tratadas.

\subsubsection{Efectos sobre precios y márgenes}

\begin{table}[H]
\centering
\caption{Efecto de la entrada sobre el precio, según el grupo de comparación}
\label{tab:control}
\tabla{tab_att_control}
\end{table}

Las tres definiciones de control entregan básicamente el mismo resultado. Para elegir el grupo de control hay que considerar dos criterios: el número de estaciones de control y la ausencia de spillovers, esto es, que ninguna estación del grupo de comparación esté ella misma expuesta al efecto de la entrada. El segundo criterio importa porque, como muestra \textcite{BUTTS2023}, cuando el control está parcialmente expuesto al tratamiento el estimador de diferencias en diferencias deja de recuperar el efecto de interés; en este caso la exposición del control tiene el mismo signo que el efecto sobre las tratadas, de modo que el sesgo atenúa la estimación hacia cero. Por lo que tanto controles lejanos como nunca-tratadas sirven, se elige nunca-tratadas. El Anexo~\ref{anexo:control} desarrolla esta comparación con el balance de covariables y los tests de tendencias previas de cada grupo.

La \Cref{fig:es} presenta el estudio de eventos de \eqref{eq:es} bajo la
especificación principal. Las trayectorias son planas antes de la entrada y caen de manera abrupta en el bin de la entrada, sin revertirse en los cinco bins posteriores, es decir, a lo largo de dos años y medio. La persistencia importa, un efecto que se disipara sería compatible con una reacción transitoria de posicionamiento, mientras que uno persistente indica un cambio en el equilibrio del mercado local, que es lo que predice el \Cref{cor:entrada}.

\begin{figure}[H]
\centering
\includegraphics[width=\textwidth]{es_control_nunca.pdf}
\caption{Efecto de la entrada sobre el precio de las incumbentes}
\label{fig:es}
\end{figure}
La \Cref{fig:ff} presenta el efecto sobre el margen para los bienes disponibles en las bases del MEPCO, con las dos definiciones de mercado local como series. Las dos filas son el mismo objeto en unidades distintas: la superior expresa el efecto en pesos por litro y la inferior como porcentaje del margen previo, de modo que la primera indica cuánto cae el margen y la segunda qué fracción de él se pierde.

\begin{figure}[H]
\centering
\includegraphics[width=\textwidth]{formato_fischer.pdf}
\caption{Efecto de la entrada sobre el margen bruto, por combustible}
\label{fig:ff}
\end{figure}

En el caso de la gasolina de 97 octanos se exhibe un efecto sistemáticamente menor y menos preciso, a un kilómetro no alcanza significancia convencional en ninguna de las dos unidades, y solo la logra al ampliar el radio a dos. La lectura que se desprende de esto es que a diferencia de la gasolina de 93 octanos la 97 es el combustible premium del mercado. Representa apenas el $3{,}3\%$ de las ventas nacionales de combustibles líquidos frente al $16\%$ de la 93 \parencite{FNE2025informe}, y sus consumidores plausiblemente presentan una demanda menos elástica al precio. Si la demanda que enfrenta la incumbente en ese producto es menos sensible, la llegada de un competidor genera menos presión para ajustarlo. El resultado es entonces informativo, pero conviene no apoyar en él ninguna conclusión que los otros dos combustibles no sostengan por sí solos.

La \Cref{hip:distancia} predice que el efecto decrece con la distancia entre la entrante y la incumbente y se anula a partir de un umbral. La
\Cref{fig:atenuacion} lo contrasta asignando cada estación al anillo de un kilómetro correspondiente a su primera entrada cercana y estimando el efecto por anillo.

\subsubsection{Atenuación por distancia}

\begin{figure}[H]
\centering
\includegraphics[width=\textwidth]{atenuacion_distancia.pdf}
\caption{Atenuación del efecto de la entrada con la distancia}
\label{fig:atenuacion}
\end{figure}

El perfil decae desde el anillo más cercano y deja de ser distinguible de cero entre los tres y los cuatro kilómetros. Por un lado, el resultado valida la predicción del modelo: como se explica en el Anexo~\ref{anexo:modelo}, lo que el perfil recupera es la proporción de incumbentes para las cuales la entrante cae dentro del radio en que la sustitución es efectiva, proporción que se agota a la distancia en que ya casi ninguna estación comparte mercado con la entrante. El umbral así estimado coincide con la delimitación de mercado geográfico que la \textcite{FNE2021informe} emplea en sus investigaciones. Por otro, es informativo a la hora de decidir utilizar o no controles dentro de los anillos entre 3 a 5km. Si el efecto sigue vivo a dos o tres kilómetros, el anillo de control necesariamente contiene estaciones parcialmente tratadas, lo que explica
por qué su estimación atenúa.

\subsubsection{Heterogeneidad según el margen previo}

\textcite{KOH2022} plantean que la respuesta a la competencia depende del margen con que la incumbente operaba antes de la entrada, una estación con márgenes holgados puede bajar precios, mientras que una que ya opera con márgenes estrechos carece de espacio para hacerlo y responde por otras vías, como reducir su probabilidad de seguir operando. Los autores no pueden contrastar esta hipótesis directamente, ya que no observan el precio mayorista y, por lo tanto, tampoco los márgenes de las estaciones, deben aproximarlos distinguiendo entre áreas metropolitanas y no metropolitanas, limitación que ellos mismos reconocen.

La disponibilidad del MEPCO permite contrastarla de manera directa. La \Cref{hip:margen} formula la misma predicción a partir de la \Cref{prop:margen}, donde el margen previo es una transformación exacta del número de competidoras efectivas de la incumbente. Para contrastarla hace falta una medida del margen con que cada estación operaba antes de la entrada, y esa medida debe cumplir dos condiciones.

La primera es que sea comparable entre estaciones que enfrentan costos y
mercados distintos. Se construye entonces como la \emph{posición relativa} de la
estación dentro de la distribución de precios de su propia comuna en cada mes.
Dentro de una comuna y un mes el costo mayorista es idéntico para todas las
estaciones y el margen porcentual es estrictamente creciente en el precio, de
modo que el orden de los precios es exactamente el orden de los márgenes. La
posición así definida es además adimensional, lo que permite promediarla entre
combustibles.

La segunda es que no esté construida sobre la misma variable que el resultado.
La medida es entonces un \emph{promedio ponderado} de la posición de la estación
en cada combustible, con ponderadores iguales a la participación relativa de cada
uno dentro de los cuatro combustibles analizados —$19{,}6\%$ la gasolina de 93
octanos, $11{,}9\%$ la de 95, $4{,}0\%$ la de 97 y $64{,}4\%$ el diésel,
calculados a partir de las ventas nacionales de 2017 que reporta
\textcite{FNE2025informe}— y \emph{excluyendo el combustible que se analiza en
cada caso}.\footnote{Los cuatro combustibles analizados concentran el $81{,}5\%$
de las ventas nacionales de combustibles líquidos, de modo que los ponderadores
son las participaciones informadas por la FNE renormalizadas sobre ese total.} Al estimar el efecto sobre el precio del diésel, por ejemplo, el
moderador se construye solo con las tres gasolinas.

La exclusión no es un refinamiento menor. El precio observado de una estación
tiene un componente persistente, que es el que interesa, y uno transitorio. Si
el moderador incluyera el propio combustible, el tercil superior quedaría
poblado en parte por estaciones que simplemente recibieron perturbaciones
positivas en la ventana previa, y esas perturbaciones revierten: el grupo alto
caería después con o sin entrada. Excluir el combustible analizado elimina ese
canal. El Anexo~\ref{anexo:moderador} compara esta definición con dos
alternativas más simples y cuantifica cuánto del contraste proviene de él.

Finalmente, la posición se mide siempre como el promedio sobre la ventana previa
a la entrada y nunca de manera contemporánea, ya que una medición posterior
sería un resultado del tratamiento y no una característica previa de la
estación.

\begin{table}[H]
\centering
\caption{Efecto de la entrada según el margen previo de la incumbente}
\label{tab:het}
\tabla{tab_het_rank}
\end{table}

La \Cref{tab:het} y la \Cref{fig:het} confirman la predicción. El efecto sobre
las estaciones del tercil superior de su comuna duplica o triplica al de las
restantes en los cinco resultados, y la fila de diferencia —que proviene de
reparametrizar el mismo modelo, de modo que su error estándar es el del
contraste y no el de una comparación visual de intervalos— es significativa al
$1\%$ en el precio del diésel y en su margen, y al $10\%$ en el precio de la
gasolina de 95 octanos. En los demás combustibles el contraste mantiene el signo
y una magnitud comparable sin alcanzar significancia convencional, de modo que en
ellos el resultado debe leerse como indicativo. Las trayectorias previas de ambos
grupos son planas, lo que descarta que el resultado sea reversión a la media.

\begin{figure}[H]
\centering
\includegraphics[width=\textwidth]{het_rank_precio.pdf}
\caption{Efecto de la entrada según el margen previo de la incumbente}
\label{fig:het}
\end{figure}

El resultado tiene una lectura que va más allá de la heterogeneidad. Lo que disciplina el precio no es la mera presencia de competidores, sino el margen que la incumbente venía capturando, que es la medida de su poder de mercado local. Es la misma conclusión a la que llega \textcite{BRESNAHAN1991} por otra vía: al inferir la conducta competitiva desde la relación entre número de firmas y tamaño del mercado en mercados locales aislados, encuentra que casi toda la variación ocurre con la entrada de la segunda o la tercera firma y que, una vez que el mercado tiene entre tres y cinco, la siguiente entrante apenas la altera. El efecto de una entrada no es constante, depende de cuán concentrado estaba el mercado que la recibe.

\subsection{Validez del supuesto de identificación}

El supuesto que sostiene el diseño no es que la entrada sea exógena, sino que lo sea el momento en que se materializa, condicional en los efectos fijos de \eqref{eq:es}. Esta subsección reúne la evidencia sobre esa afirmación. Los contrastes de tendencias previas, que son la verificación indirecta habitual, se presentan en el Anexo~\ref{anexo:control} junto con la comparación de grupos de control.

Conviene delimitar la afirmación antes de someterla a prueba, porque la mitad de las objeciones habituales apuntan a algo que el diseño no necesita. Como las cadenas eligen deliberadamente dónde instalarse, resulta inverosímil que la ubicación sea independiente de las condiciones del mercado local. La teoría lo anticipa desde el origen mismo de este marco: \textcite{HOTELLING1929} muestra que la competencia espacial genera un incentivo a instalarse cerca del rival y no lejos de él, y la evidencia estructural para este mismo mercado confirma que entrar y salir son decisiones dinámicas que responden a la rentabilidad esperada del mercado local y al marco regulatorio de precios que lo rige \parencite{HOUDE2005}. Los datos de este trabajo son consistentes con ambas cosas: el $72{,}6\%$ de las entradas ocurre teniendo al menos una estación vecina dentro de 2 kilómetros. Lo que se supone, entonces, no es que la ubicación sea exógena, sino que, condicional en el efecto fijo de estación y en los efectos de tiempo, la fecha precisa de apertura no responde a la evolución de los precios de las incumbentes.

Esta separación tiene un fundamento institucional en el caso chileno. La decisión de instalar una estación es endógena, pero entre esa decisión y la apertura media un rezago largo, variable y de naturaleza administrativa: adquisición o arriendo del terreno, permiso de edificación, exigencias del plan regulador comunal, autorización de la Superintendencia de Electricidad y Combustibles y recepción de obras. La \textcite{FNE2021informe} documenta que la escasez de terrenos aptos, precisamente por las exigencias de los planos reguladores en las zonas urbanas más saturadas, constituye la principal barrera a la entrada del mercado. Es ese rezago el que hace que el momento de apertura sea, desde la perspectiva de la incumbente, un evento cuya fecha no está ligada a su propia conducta de precios.

La \Cref{tab:escalera} estima \eqref{eq:did} sobre la misma muestra y con el mismo regresor, variando únicamente los efectos fijos. La comparación es la de la tabla 6 de \textcite{ARCIDIACONO2020}.

\begin{table}[H]
\centering
\caption{Efecto estimado según los efectos fijos incluidos}
\label{tab:escalera}
\tabla{tab_fe_ladder}
\end{table}

La especificación de sección cruzada arroja un efecto cercano al triple del que entrega la especificación preferida. Esa brecha es la medida de cuánto de la asociación bruta entre entrada y precios bajos proviene de dónde eligen instalarse las entrantes, y no de lo que su entrada provoca. El efecto fijo de estación es el que la elimina.

